%!TEX root = ../../main.tex
% PPT绘图说明：画五条水平泳道，分别对应Qwen-RobotManip、Qwen-RobotWorld、pi family、GR00T和Gemini Robotics。
% 每条泳道统一使用source roles、preparation operations、training stages、feedback/validation四列。
% 用相同颜色表示相同数据角色，重点突出系统之间可比较的结构，不放不可直接比较的规模数字或模型分数。
% 图的最下方总结共同骨架：align sources -> assign roles -> train -> evaluate -> update data。
\begin{figure*}[t]
\centering
\fbox{\rule{0pt}{6.4cm}\rule{0.96\textwidth}{0pt}}
\caption{Planned structural comparison of public industrial data recipes. The final figure should align systems by source role, preparation operation, training stage, and feedback rather than by incomparable data volumes or end-to-end scores.}
\label{fig:industrial-recipe-map}
\end{figure*}
